\chapter{Angles et cercle trigonométrique}
\label{manual:chapter:14}

\begin{questionbox}
Comment un angle devient-il un nombre, et comment le cercle unité permet-il de définir sinus et cosinus pour tous les angles?
\end{questionbox}

\section{Degrés et radians}
\label{manual:section:14.01}

Un degré partage un tour complet en $360$ parties. Le radian relie directement un angle à une longueur d'arc.

\begin{definition}{Mesure en radians}{}
Pour un cercle de rayon $r$, un angle au centre qui intercepte un arc de longueur $s$ possède la mesure
\[
\theta=\frac sr
\]
en radians.
\end{definition}

Un tour complet a une longueur $2\pi r$, donc mesure $2\pi$ radians. Ainsi,
\[
180^\circ=\pi\text{ radians}.
\]
Les conversions sont
\[
\theta_{\mathrm{rad}}=\theta_{\mathrm{deg}}\frac{\pi}{180},
\qquad
\theta_{\mathrm{deg}}=\theta_{\mathrm{rad}}\frac{180}{\pi}.
\]

\begin{examplebox}
\[
60^\circ=60\frac{\pi}{180}=\frac\pi3,
\qquad
\frac{5\pi}{6}=\frac{5\pi}{6}\frac{180}{\pi}=150^\circ.
\]
\end{examplebox}

\section{Le cercle unité}
\label{manual:section:14.02}

\begin{definition}{Cercle trigonométrique}{}
Le cercle trigonométrique est le cercle de centre $(0,0)$ et de rayon $1$. Le point associé à l'angle $\theta$, mesuré depuis l'axe positif des $x$ dans le sens antihoraire, est
\[
P(\cos\theta,\sin\theta).
\]
\end{definition}

\begin{center}
\begin{tikzpicture}[scale=2.3]
\draw[-{Stealth},thick] (-1.25,0)--(1.35,0) node[right] {$x$};
\draw[-{Stealth},thick] (0,-1.25)--(0,1.35) node[above] {$y$};
\draw[thick,uqamblue] (0,0) circle (1);
\def\ang{42}
\coordinate (P) at ({cos(\ang)},{sin(\ang)});
\draw[very thick,teal,-{Stealth[length=2mm]}] (0,0)--(P);
\draw[dashed] (P)--({cos(\ang)},0);
\draw[dashed] (P)--(0,{sin(\ang)});
\draw[gold,thick] (0.30,0) arc[start angle=0,end angle=42,radius=0.30];
\node[gold] at (0.42,0.16) {$\theta$};
\fill[deepblue] (P) circle (0.025) node[above right] {$P(\cos\theta,\sin\theta)$};
\node[below] at ({cos(\ang)/2},0) {$\cos\theta$};
\node[left] at (0,{sin(\ang)/2}) {$\sin\theta$};
\end{tikzpicture}
\end{center}

Comme $P$ est sur le cercle $x^2+y^2=1$, on obtient immédiatement l'identité fondamentale.

\begin{theorem}{Identité pythagoricienne}{}
Pour tout angle réel $\theta$,
\[
\cos^2\theta+\sin^2\theta=1.
\]
\end{theorem}

\section{Angles remarquables}
\label{manual:section:14.03}

Les valeurs exactes proviennent des triangles $45^\circ$-$45^\circ$-$90^\circ$ et $30^\circ$-$60^\circ$-$90^\circ$.

\begin{center}
\begin{tabular}{c|ccccc}
\toprule
$\theta$ & $0$ & $\pi/6$ & $\pi/4$ & $\pi/3$ & $\pi/2$ \\
\midrule
$\sin\theta$ & $0$ & $1/2$ & $\sqrt2/2$ & $\sqrt3/2$ & $1$ \\
$\cos\theta$ & $1$ & $\sqrt3/2$ & $\sqrt2/2$ & $1/2$ & $0$ \\
$\tan\theta$ & $0$ & $1/\sqrt3$ & $1$ & $\sqrt3$ & non définie \\
\bottomrule
\end{tabular}
\end{center}

Dans les autres quadrants, les valeurs absolues se retrouvent avec l'angle de référence; les signes proviennent des coordonnées du point sur le cercle.

\begin{examplebox}
L'angle $150^\circ$ a pour angle de référence $30^\circ$ et se trouve dans le deuxième quadrant. Ainsi,
\[
\sin150^\circ=\frac12,
\qquad
\cos150^\circ=-\frac{\sqrt3}{2},
\qquad
\tan150^\circ=-\frac1{\sqrt3}.
\]
\end{examplebox}

\section{Périodicité et symétries}
\label{manual:section:14.04}

Un tour complet ramène au même point :
\[
\sin(\theta+2\pi)=\sin\theta,
\qquad
\cos(\theta+2\pi)=\cos\theta.
\]
La tangente, définie par
\[
\tan\theta=\frac{\sin\theta}{\cos\theta},
\]
a une période $\pi$ et n'est pas définie lorsque $\cos\theta=0$.

Les symétries du cercle donnent
\[
\sin(-\theta)=-\sin\theta,
\qquad
\cos(-\theta)=\cos\theta.
\]
Le sinus est impair et le cosinus est pair.


\section{Le radian mesure naturellement une rotation}
\label{manual:section:14.05}

\begin{visualbox}
Un degré dépend d'une convention de partage du cercle en $360$ parties. Un radian compare deux longueurs : l'arc parcouru et le rayon. Cette définition explique pourquoi les formules trigonométriques utilisent naturellement $\pi$ et pourquoi, sur le cercle unité, la mesure de l'angle est exactement la longueur de l'arc.
\end{visualbox}

\begin{center}
\begin{tikzpicture}[scale=2.0]
  \draw[slate,thick] (0,0) circle (1);
  \draw[navy,very thick] (0,0)--(1,0);
  \draw[navy,very thick] (0,0)--({cos(57.3)},{sin(57.3)});
  \draw[teal,very thick] (1,0) arc[start angle=0,end angle=57.3,radius=1];
  \node[teal] at (0.85,0.48) {$s=r$};
  \node[navy] at (0.5,-0.13) {$r$};
  \node[navy] at (0.28,0.52) {$1\text{ rad}$};
\end{tikzpicture}
\qquad
\begin{tikzpicture}[scale=2.0]
  \draw[slate,thick] (0,0) circle (1);
  \draw[-{Stealth},slate] (-1.25,0)--(1.3,0);
  \draw[-{Stealth},slate] (0,-1.25)--(0,1.3);
  \coordinate (P) at ({cos(140)},{sin(140)});
  \draw[gold,very thick] (0,0)--(P);
  \draw[dashed,teal] (P)--({cos(140)},0);
  \fill[terracotta] (P) circle (2pt);
  \node[anchor=south east] at (P) {$P(\cos\theta,\sin\theta)$};
\end{tikzpicture}
\end{center}

\begin{connectionbox}
Sur le cercle unité, le cosinus est la projection horizontale du rayon et le sinus sa projection verticale. L'identité
\[
\cos^2\theta+\sin^2\theta=1
\]
n'est donc pas une formule supplémentaire : c'est simplement l'équation $x^2+y^2=1$ du cercle appliquée au point $P$.
\end{connectionbox}

\subsection{Angles de référence et symétries}
Plutôt que mémoriser une longue table, on localise le quadrant, on ramène l'angle à un angle aigu de référence, puis on détermine les signes des coordonnées. Les symétries du cercle expliquent les relations comme $\cos(-\theta)=\cos\theta$ et $\sin(-\theta)=-\sin\theta$.


\section{Pourquoi les radians simplifient les formules}
\label{manual:section:14.06}

Pour un angle $\theta$ en radians, la longueur d'arc est
\[
s=r\theta.
\]
Cette formule ne contient aucun facteur de conversion parce que le radian est précisément défini par le rapport $s/r$. En degrés, il faudrait écrire
\[
s=r\theta\frac{\pi}{180}.
\]
Le radian est donc l'unité qui rend la relation entre angle et longueur directement proportionnelle.

\subsection{Une petite rotation produit un petit déplacement}
Sur un cercle de rayon \(r>0\), une variation angulaire \(\Delta\theta\) mesurée en radians produit un déplacement \emph{orienté} le long du cercle exactement égal à
\[
\Delta s=r\Delta\theta.
\]
La longueur parcourue est \(r|\Delta\theta|\) pour une rotation sans changement de sens. Ce n'est pas la longueur de la corde reliant les deux positions.
Cette idée est la base de la vitesse angulaire et explique pourquoi les modèles de mouvement circulaire utilisent les radians.

\begin{center}
\begin{tikzpicture}[scale=1.5]
 \draw[slate,thick] (0,0) circle (2);
 \draw[navy,thick] (0,0)--(2,0);
 \draw[navy,thick] (0,0)--({2*cos(25)},{2*sin(25)});
 \draw[teal,very thick] (2,0) arc[start angle=0,end angle=25,radius=2];
 \draw[gold,thick,<->] (0.65,0) arc[start angle=0,end angle=25,radius=0.65];
 \node[gold!75!black] at (0.78,0.2) {$\Delta\theta$};
 \node[teal] at (1.65,0.45) {$\Delta s$};
 \node[navy] at (1,-0.15) {$r$};
\end{tikzpicture}
\end{center}

\subsection{Construire les valeurs remarquables}
Les valeurs de $30^\circ$, $45^\circ$ et $60^\circ$ ne doivent pas être apprises comme une liste sans origine. Elles viennent de deux triangles : un triangle équilatéral coupé en deux et un triangle rectangle isocèle.

\begin{center}
\begin{tikzpicture}[scale=0.95]
 \draw[navy,very thick] (0,0)--(4,0)--(2,{2*sqrt(3)})--cycle;
 \draw[teal,thick,dashed] (2,0)--(2,{2*sqrt(3)});
 \node[below] at (1,0) {$1$}; \node[below] at (3,0) {$1$};
 \node[left] at (1,{sqrt(3)}) {$2$};
 \node[right] at (2.1,{sqrt(3)}) {$\sqrt3$};
 \node[gold!75!black] at (0.65,0.28) {$60^\circ$};
 \node[gold!75!black] at (2.25,0.35) {$30^\circ$};
\end{tikzpicture}
\qquad
\begin{tikzpicture}[scale=0.95]
 \draw[navy,very thick] (0,0)--(3,0)--(3,3)--cycle;
 \draw (2.7,0) rectangle (3,0.3);
 \node[below] at (1.5,0) {$1$}; \node[right] at (3,1.5) {$1$};
 \node[above left] at (1.5,1.5) {$\sqrt2$};
 \node[gold!75!black] at (0.75,0.3) {$45^\circ$};
\end{tikzpicture}
\end{center}

Ces constructions rendent les coordonnées du cercle unité prévisibles et réduisent la mémorisation à deux figures fondamentales.


\section{Mesurer et lire une rotation}
\label{manual:section:14.07}

Les degrés partagent un tour en $360$ parts; les radians mesurent directement une longueur d'arc rapportée au rayon. Les deux unités décrivent la même rotation, mais le radian relie naturellement angle, arc, vitesse angulaire et fonctions trigonométriques.

\subsection{Le radian comme rapport de longueurs}

Dans un cercle de rayon $r$, un angle central $\theta$ radians intercepte un arc de longueur
\[
s=r\theta.
\]
Ainsi
\[
\theta=\frac sr.
\]
Le rapport est sans unité, car le numérateur et le dénominateur sont deux longueurs. Sur le cercle unité, $r=1$, donc la mesure de l'angle est numériquement égale à la longueur de l'arc.

Un tour complet possède une longueur $2\pi r$, donc
\[
2\pi\ \text{radians}=360^\circ.
\]
On en déduit
\[
180^\circ=\pi,
\qquad
1^\circ=\frac\pi{180}.
\]

\subsection{Conversions raisonnées}

Pour convertir des degrés en radians, on multiplie par $\pi/180$ :
\[
150^\circ=150\cdot\frac\pi{180}=\frac{5\pi}{6}.
\]
Pour convertir des radians en degrés, on multiplie par $180/\pi$ :
\[
\frac{7\pi}{4}\cdot\frac{180}{\pi}=315^\circ.
\]
Les facteurs de conversion sont égaux à $1$ puisque $\pi$ radians et $180^\circ$ désignent le même angle.

\subsection{Le cercle unité comme système de coordonnées}

À l'angle $\theta$ correspond le point
\[
P(\theta)=(\cos\theta,\sin\theta).
\]
Le cosinus est l'abscisse et le sinus l'ordonnée. Comme $P$ appartient au cercle de rayon $1$,
\[
\cos^2\theta+\sin^2\theta=1.
\]
Cette identité fondamentale est simplement l'équation du cercle
\[
x^2+y^2=1
\]
appliquée aux coordonnées de $P$.

\subsection{Construire les valeurs remarquables}

Les valeurs de $45^\circ$ proviennent d'un triangle rectangle isocèle. Si les deux côtés de l'angle droit valent $1$, l'hypoténuse vaut $\sqrt2$. Après normalisation sur le cercle unité,
\[
\cos45^\circ=\sin45^\circ=\frac1{\sqrt2}=\frac{\sqrt2}{2}.
\]

Pour $30^\circ$ et $60^\circ$, on coupe un triangle équilatéral de côté $2$ en deux. On obtient un triangle rectangle de côtés
\[
1,\quad\sqrt3,\quad2.
\]
Ainsi
\[
\sin30^\circ=\frac12,
\qquad
\cos30^\circ=\frac{\sqrt3}{2},
\]
\[
\sin60^\circ=\frac{\sqrt3}{2},
\qquad
\cos60^\circ=\frac12.
\]
Les valeurs ne sont donc pas une table arbitraire à mémoriser; elles viennent de deux triangles fondamentaux.

\subsection{Angles coterminaux}

Deux angles qui diffèrent d'un nombre entier de tours aboutissent au même point :
\[
\theta+2k\pi,
\qquad k\in\Z.
\]
Ainsi
\[
\cos(\theta+2k\pi)=\cos\theta,
\qquad
\sin(\theta+2k\pi)=\sin\theta.
\]
Pour ramener un angle au tour principal, on ajoute ou soustrait des multiples de $2\pi$.

\begin{examplebox}
L'angle
\[
\frac{17\pi}{6}
\]
diffère de $2\pi=12\pi/6$ de l'angle
\[
\frac{5\pi}{6}.
\]
Les deux ont donc les mêmes valeurs trigonométriques.
\end{examplebox}

\subsection{Angle de référence et signes par quadrant}

L'angle de référence est l'angle aigu formé avec l'axe horizontal. Il permet de réutiliser les valeurs remarquables, puis de déterminer les signes à partir du quadrant :
\begin{center}
\begin{tabular}{c|ccc}
\toprule
Quadrant & $\cos$ & $\sin$ & $\tan$\\\midrule
I & $+$ & $+$ & $+$\\
II & $-$ & $+$ & $-$\\
III & $-$ & $-$ & $+$\\
IV & $+$ & $-$ & $-$\\
\bottomrule
\end{tabular}
\end{center}

Par exemple, $330^\circ$ a pour angle de référence $30^\circ$ et se trouve dans le quatrième quadrant :
\[
\cos330^\circ=\frac{\sqrt3}{2},
\qquad
\sin330^\circ=-\frac12,
\qquad
\tan330^\circ=-\frac1{\sqrt3}.
\]

\subsection{Symétries}

La réflexion par rapport à l'axe horizontal donne
\[
\cos(-\theta)=\cos\theta,
\qquad
\sin(-\theta)=-\sin\theta.
\]
Le cosinus est pair et le sinus impair. Une rotation d'un demi-tour donne
\[
\cos(\theta+\pi)=-\cos\theta,
\qquad
\sin(\theta+\pi)=-\sin\theta.
\]
Ces relations se voient directement sur les coordonnées du cercle.

\subsection{Longueur d'arc et aire d'un secteur}

Lorsque $\theta$ est en radians,
\[
s=r\theta
\]
pour la longueur d'arc, et
\[
A=\frac12r^2\theta
\]
pour l'aire du secteur. La seconde formule vient du rapport
\[
\frac{\theta}{2\pi}
\]
entre l'angle et le tour complet :
\[
A=\frac{\theta}{2\pi}\cdot\pi r^2=\frac12r^2\theta.
\]

\subsection{Degrés ou radians sur la calculatrice}

Une calculatrice en mode degré interprète $\sin(1)$ comme $\sin(1^\circ)$; en mode radian, elle interprète $\sin(1)$ comme $\sin(1\ \text{rad})$. Les valeurs sont très différentes. Avant tout calcul :
\begin{itemize}
\item repérer le symbole $^\circ$ ou les multiples de $\pi$;
\item choisir le mode correspondant;
\item tester une valeur connue, par exemple $\sin30^\circ=1/2$ ou $\sin(\pi/6)=1/2$.
\end{itemize}

\begin{summarybox}
Le radian mesure une rotation par un rapport de longueurs. Le cercle unité transforme l'angle en coordonnées, explique l'identité pythagoricienne, les signes, les symétries et la périodicité des fonctions trigonométriques.
\end{summarybox}


